% GENERATED FILE. Edit canonical lessons, then run npm run generate:books.
% canonical-source-hash: fnv1a64:59038787e2f19e61
% canonical-lessons: LA-C136-cubitum, LA-C136-pollex, LA-C136-cutis, LA-C136-os, LA-C136-stomachus

\chapter{Elbow, Thumb, Skin, Bone, Stomach}
\label{ch:la-elbow-thumb-skin-bone-stomach}
\hlchaptermodality{136}

\begin{chapteropening}
\textbf{By the end of this chapter:} \emph{I can say an elbow, a thumb, the skin, a bone and the stomach.}

Complete the last lesson of chapter 136: I can say an elbow, a thumb, the skin, a bone and the stomach.
\end{chapteropening}

\section[cubitum, cubitī]{cubitum, cubitī — an elbow}

\label{lesson:LA-C136-cubitum}



\begin{quote}
Before the new one: say the Latin for the chin, then the Latin for the back.
\end{quote}

\subsection*{You'll want to know: cubitum, cubitī}
\textbf{cubitum, cubitī} — \textquotedblleft{}an elbow\textquotedblright{}.

\begin{cousinweb}
\textbf{cubitum} is built on \emph{cubō}, \textquotedblleft{}to lie down, to recline\textquotedblright{} (EtymDB 2.1): the elbow is what one leans on at table. English \emph{cubit}, the length from elbow to fingertip, comes from it.
\end{cousinweb}

\begin{grammarlens}[title={What you've built}]
One more word for this chapter.
\end{grammarlens}

\subsection*{Your turn}
\begin{itemize}
\raggedright
  \item \emph{Say it:} \emph{cubitum}
  \item \emph{Say it:} \emph{cubitum}, once more
  \item \emph{Say it:} say \emph{cubitum}
  \item \emph{Recall:} say the Latin for the chin, then the Latin for the back, then say \emph{cubitum} again
\end{itemize}

\begin{tcolorbox}[breakable,colback=teal!4,colframe=teal!35!black,title={Before you move on}]
What does cubitum, cubitī mean? (\textquotedblleft{}An elbow\textquotedblright{}.) Say it once more.
\end{tcolorbox}
\section[pollex, pollicis]{pollex, pollicis — a thumb}

\label{lesson:LA-C136-pollex}



\begin{quote}
Before the new one: say the Latin for the back, then the Latin for an elbow.
\end{quote}

\subsection*{You'll want to know: pollex, pollicis}
\textbf{pollex, pollicis} — \textquotedblleft{}a thumb\textquotedblright{}.

\begin{cousinweb}
\textbf{pollex} was linked by the Romans with \emph{polleō}, \textquotedblleft{}to be strong\textquotedblright{} (Lewis \& Short): the strongest finger. Modern scholars are less sure.
\end{cousinweb}

\begin{grammarlens}[title={What you've built}]
One more word for this chapter.
\end{grammarlens}

\subsection*{Your turn}
\begin{itemize}
\raggedright
  \item \emph{Say it:} \emph{pollex}
  \item \emph{Say it:} \emph{pollex}, once more
  \item \emph{Say it:} say \emph{pollex}
  \item \emph{Recall:} say the Latin for the back, then the Latin for an elbow, then say \emph{pollex} again
\end{itemize}

\begin{tcolorbox}[breakable,colback=teal!4,colframe=teal!35!black,title={Before you move on}]
What does pollex, pollicis mean? (\textquotedblleft{}A thumb\textquotedblright{}.) Say it once more.
\end{tcolorbox}
\section[cutis, cutis]{cutis, cutis — the skin}

\label{lesson:LA-C136-cutis}



\begin{quote}
Before the new one: say the Latin for an elbow, then the Latin for a thumb.
\end{quote}

\subsection*{You'll want to know: cutis, cutis}
\textbf{cutis, cutis} — \textquotedblleft{}the skin\textquotedblright{}.

\begin{cousinweb}
\textbf{cutis} is inherited, usually traced to Proto-Indo-European \emph{*kuHt-}, \textquotedblleft{}skin, hide\textquotedblright{} (EtymDB 2.1); Greek \emph{kytos}, \textquotedblleft{}a hollow, a skin\textquotedblright{}, is compared (Lewis \& Short), and English \emph{hide} may be a cousin.
\end{cousinweb}

\begin{grammarlens}[title={What you've built}]
One more word for this chapter.
\end{grammarlens}

\subsection*{Your turn}
\begin{itemize}
\raggedright
  \item \emph{Say it:} \emph{cutis}
  \item \emph{Say it:} \emph{cutis}, once more
  \item \emph{Say it:} say \emph{cutis}
  \item \emph{Recall:} say the Latin for an elbow, then the Latin for a thumb, then say \emph{cutis} again
\end{itemize}

\begin{tcolorbox}[breakable,colback=teal!4,colframe=teal!35!black,title={Before you move on}]
What does cutis, cutis mean? (\textquotedblleft{}The skin\textquotedblright{}.) Say it once more.
\end{tcolorbox}
\section[os, ossis]{os, ossis — a bone}

\label{lesson:LA-C136-os}



\begin{quote}
Before the new one: say the Latin for a thumb, then the Latin for the skin.
\end{quote}

\subsection*{You'll want to know: os, ossis}
\textbf{os, ossis} — \textquotedblleft{}a bone\textquotedblright{}.

\begin{cousinweb}
\textbf{os}, \emph{ossis}, is inherited, usually traced to Proto-Indo-European \emph{*h3esth1}, \textquotedblleft{}bone\textquotedblright{} (EtymDB 2.1); Greek \emph{osteon}, \textquotedblleft{}bone\textquotedblright{}, is a cousin (Lewis \& Short). It is a different word from \emph{ōs}, \emph{ōris}, \textquotedblleft{}mouth\textquotedblright{}.
\end{cousinweb}

\begin{grammarlens}[title={What you've built}]
One more word for this chapter.
\end{grammarlens}

\subsection*{Your turn}
\begin{itemize}
\raggedright
  \item \emph{Say it:} \emph{os}
  \item \emph{Say it:} \emph{os}, once more
  \item \emph{Say it:} say \emph{os}
  \item \emph{Recall:} say the Latin for a thumb, then the Latin for the skin, then say \emph{os} again
\end{itemize}

\begin{tcolorbox}[breakable,colback=teal!4,colframe=teal!35!black,title={Before you move on}]
What does os, ossis mean? (\textquotedblleft{}A bone\textquotedblright{}.) Say it once more.
\end{tcolorbox}
\section[stomachus, stomachī]{stomachus, stomachī — the stomach}

\label{lesson:LA-C136-stomachus}



\begin{quote}
Before the new one: say the Latin for the skin, then the Latin for a bone.
\end{quote}

\subsection*{You'll want to know: stomachus, stomachī}
\textbf{stomachus, stomachī} — \textquotedblleft{}the stomach\textquotedblright{}.

\begin{cousinweb}
\textbf{stomachus} was borrowed from Greek \emph{stomakhos}, \textquotedblleft{}throat, gullet\textquotedblright{} (Lewis \& Short), from \emph{stoma}, \textquotedblleft{}mouth\textquotedblright{} (EtymDB 2.1).
\end{cousinweb}

\begin{grammarlens}[title={What you've built}]
One more word for this chapter.
\end{grammarlens}

\subsection*{Your turn}
\begin{itemize}
\raggedright
  \item \emph{Say it:} \emph{stomachus}
  \item \emph{Say it:} \emph{stomachus}, once more
  \item \emph{Say it:} say \emph{stomachus}
  \item \emph{Recall:} say the Latin for the skin, then the Latin for a bone, then say \emph{stomachus} again
\end{itemize}

\begin{tcolorbox}[breakable,colback=teal!4,colframe=teal!35!black,title={Before you move on}]
What does stomachus, stomachī mean? (\textquotedblleft{}The stomach\textquotedblright{}.) Say it once more.
\end{tcolorbox}
